\documentclass[../../main.tex]{subfiles}

\begin{document}

\section{积分的可数可加性与连续性}\label{section:ch04-可数可加连续}

我们在前一节证明的 Lebesgue 积分的线性与单调性质是 Riemann 积分熟知的性质的推广. 在这简短的一节中, 我们证明 Riemann 积分所没有的两个 Lebesgue 积分的性质. 以下关于 Lebesgue 积分的可数可加性是 Lebesgue 测度的可数可加性的伴随.

\begin{theorem}[积分的可数可加性]\label{theorem:ch04-可数可加}
令 $f$ 为 $E$ 上可积而 $\{E_{n}\}_{n=1}^{\infty}$ 是 $E$ 的可测子集的不交可数族, 其并集是 $E$. 则
\begin{align}
\int_{E}f=\sum_{n=1}^{\infty}\int_{E_{n}}f.\label{align:ch04-可数可加}
\end{align}
\end{theorem}

\begin{proof}
令 $n$ 为自然数. 定义 $f_{n}=f\cdot\chi_{n}$, 其中 $\chi_{n}$ 是可测集 $\bigcup_{k=1}^{n}E_{k}$ 的特征函数. 则 $f_{n}$ 是 $E$ 上的可测函数且
\begin{align*}
\text{在 }E\text{ 上},\quad|f_{n}|\leqslant|f|.
\end{align*}
观察到在 $E$ 上逐点 $\{f_{n}\}\to f$. 因此, 根据 Lebesgue 控制收敛定理,
\begin{align*}
\int_{E}f=\lim_{n\to\infty}\int_{E}f_{n}.
\end{align*}
另一方面, 由 $\{E_{n}\}_{n=1}^{\infty}$ 是不交的, 从积分在区域上的可加性得出, 对每个 $n$,
\begin{align*}
\int_{E}f_{n}=\sum_{k=1}^{n}\int_{E_{k}}f.
\end{align*}
因此
\begin{align*}
\int_{E}f=\lim_{n\to\infty}\int_{E}f_{n}=\lim_{n\to\infty}\left[\sum_{k=1}^{n}\int_{E_{k}}f\right]=\sum_{n=1}^{\infty}\int_{E_{n}}f.
\end{align*}
\end{proof}

我们把用积分的可数可加性证明关于积分的连续性的以下结果留给读者: 它可作为基于测度的可数可加性证明测度的连续性的一个模式.

\begin{theorem}[积分的连续性]\label{theorem:ch04-积分连续}
令 $f$ 在 $E$ 上可积.
\begin{enumerate}[(i)]
\item 若 $\{E_{n}\}_{n=1}^{\infty}$ 是 $E$ 的可测子集的上升的可数族, 则
\begin{align}
\int_{\bigcup_{n=1}^{\infty}E_{n}}f=\lim_{n\to\infty}\int_{E_{n}}f.\label{align:ch04-上升连续}
\end{align}
\item 若 $\{E_{n}\}_{n=1}^{\infty}$ 是 $E$ 的可测子集的下降的可数族, 则
\begin{align}
\int_{\bigcap_{n=1}^{\infty}E_{n}}f=\lim_{n\to\infty}\int_{E_{n}}f.\label{align:ch04-下降连续}
\end{align}
\end{enumerate}
\end{theorem}

\begin{exercise}\label{exercise:ch04-37}
令 $f$ 为 $E$ 上的可积函数. 证明对每个 $\varepsilon>0$, 存在自然数 $N$ 使得若 $n\geqslant N$, 则 $\left|\int_{E_{n}}f\right|<\varepsilon$, 其中 $E_{n}=\{x\in E\mid|x|\geqslant n\}$.
\end{exercise}

\begin{exercise}\label{exercise:ch04-38}
对如下定义在 $[1,\infty)$ 上的两个函数中的每个, 证明 $\lim_{n\to\infty}\int_{1}^{n}f$ 存在而 $f$ 在 $[1,\infty)$ 上不可积. 这与积分的连续性矛盾吗?
\begin{enumerate}[(i)]
\item 定义 $f(x)=(-1)^{n}/n$, $n\leqslant x<n+1$.
\item 定义 $f(x)=(\sin x)/x$, $1\leqslant x<\infty$.
\end{enumerate}
\end{exercise}

\begin{exercise}\label{exercise:ch04-39}
证明关于积分的连续性的定理.
\end{exercise}

\end{document}